% !TEX root = ../main.tex
\chapter{Eurodollars, Parallel Settlement}\label{ch:eurodollar}
\begin{paracol}{2}

\begin{sources}
\begin{tabularx}{\linewidth}{@{}l l L@{}}
\toprule
\textbf{Segment} & \textbf{Length} & \textbf{Content}\\
\midrule
\seg{8}{1}{L8.1 FT: Ring-fencing and the Volcker Rule} & 9:45 & Ring-fencing retail banking; matched book versus proprietary trading.\\
\seg{8}{2}{L8.2 The Eurodollar Market in Crisis} & 4:51 & LIBOR--OIS in 2007--2009.\\
\seg{8}{3}{L8.3 What are Eurodollars?} & 7:11 & Stigum's Exxon example; dollar deposits outside the US with reserves at a New York correspondent.\\
\seg{8}{4}{L8.4 Why is there a Eurodollar Market?} & 4:43 & Demand for dollars abroad; capital controls; tail and dog.\\
\seg{8}{5}{L8.5 Eurodollar as Global Funding Market} & 11:33 & Eurobanks as money dealers; the interbank market and LIBOR.\\
\seg{8}{6}{L8.6 Liquidity Challenge of Eurodollar Banks} & 10:59 & Lining up cash flows in time; a forward deposit as a swap of IOUs.\\
\seg{8}{7}{L8.7 FRA as Implicit Swap of IOUs} & 4:34 & Forward forwards and forward rate agreements, off balance sheet.\\
\seg{8}{8}{L8.8 Forward Parity, Interest Rates, EH} & 8:58 & Forward interest parity; covered interest parity via a swap of deposits.\\
\seg{8}{9}{L8.9 Forward Parity, Exchange Rates, UIP} & 2:48 & Forward rates are not expected spot rates: EH and UIP fail.\\
\seg{8}{10}{L8.10 Forward Rates are NOT Expected Spot Rates} & 2:49 & (Repeats L8.9.)\\
\bottomrule
\end{tabularx}

\smallskip
\textbf{Files.} Transcripts \file{materials/transcripts/L08-P*.txt}.
\textbf{Reading.} Stigum, \emph{The Money Market}, chapter on the Eurodollar market (from p.~212), and on FRAs.
\end{sources}

\section*{Motivation}
The third of the three wholesale dollar money markets, after Fed funds and repo \ts{8}{2}{0:07}. Eurodollars are dollar deposits at banks outside
the United States, which hold their dollar reserves not at the Fed but at a correspondent bank in New York: a \emph{parallel} settlement system, one
level below the Fed funds market in the hierarchy. The lecture explains what Eurodollars are and why they exist, how Eurobanks act as dealers in money
and line up their cash flows in time, how this leads to the first derivatives of the course, forward rate agreements, understood as implicit swaps of
IOUs, and to two parity conditions, one of which fails in practice.

% ------------------------------------------------------------------
\section{FT: ring-fencing and the Volcker rule}\label{sec:ed-ft}

\begin{casestudy}[FT, autumn 2012: ``Banks braced for proposals to ring-fence trading'']
A full-page FT feature, ``Leaner and meaner'', speculates on the business model of banking after the crisis: banks have learned some lessons, and the
regulatory landscape has yet to firm up \ts{8}{1}{0:07}\ts{8}{1}{0:28}. The article on page five refers to two ideas: the UK Vickers Commission's
recommendation that banks' retail operations be \emph{ring-fenced}, and the US \emph{Volcker rule} limiting proprietary trading \ts{8}{1}{0:48}\ts{8}{1}{1:08}.
\end{casestudy}

\paragraph{Ring-fencing.} Take the retail side of a bank, retail deposits and retail loans (households, and businesses: the distinction is retail versus
\emph{wholesale}, such as Eurodollar borrowing or derivatives), allocate to it part of the bank's capital and reserves, and fence it off. Everything else,
trading assets, wholesale borrowing, must have its own capital and reserves, so that if it loses a lot of money it can fail while the retail part is carved
off and kept safe \ts{8}{1}{1:29}\ts{8}{1}{1:50}\ts{8}{1}{2:14}\ts{8}{1}{2:34}\ts{8}{1}{2:56}\ts{8}{1}{3:19}\ts{8}{1}{3:39}.

\paragraph{The Volcker rule.} A bank may not do \emph{proprietary trading}, i.e.\ speculate on the direction of a price; it may do \emph{market making},
a matched book \ts{8}{1}{4:20}\ts{8}{1}{4:41}. Example: a pension fund holds a risky asset and wants to hedge it with a \emph{credit default swap}, a kind of
insurance, though not regulated as insurance and not a large-numbers game, that pays it if the asset fails \ts{8}{1}{5:01}\ts{8}{1}{5:22}\ts{8}{1}{5:45}\ts{8}{1}{6:06}.
The bank may sell it, but then must find another client and buy an offsetting CDS: on both sides, earning the spread, like the repo dealer of lecture~7 but with a
derivative \ts{8}{1}{6:26}\ts{8}{1}{6:47}\ts{8}{1}{7:08}. What is outlawed is the naked position, betting that the assets will fail \ts{8}{1}{7:28}\ts{8}{1}{7:48}\ts{8}{1}{8:09}.
Must the hedge be the same instrument? What matters is the same risk exposure; but hedging with something only correlated leaves \emph{basis risk}: is that
proprietary trading? If it is outlawed, banks cannot hedge, and cannot make the market at all. ``The devil is in the details'' \ts{8}{1}{8:30}\ts{8}{1}{8:51}\ts{8}{1}{9:12}\ts{8}{1}{9:32}.
\begin{taccts}
\tacct[2.0cm]{Client 1 (pension fund)}{Risky asset\\CDS bought}{}\quad
\tacct[2.0cm]{Bank (matched book)}{CDS bought from client 2}{CDS sold to client 1}\quad
\tacct[2.0cm]{Client 2}{}{CDS sold}
\end{taccts}\begin{history}[Vickers and Volcker]
The UK Independent Commission on Banking, chaired by Sir John Vickers, published its final report in September 2011, recommending ring-fencing of retail
banking; it was enacted in the Financial Services (Banking Reform) Act 2013 and in force from 2019. The Volcker rule, proposed by former Fed chairman Paul
Volcker, is section 619 of the Dodd--Frank Act (2010); the final implementing rule was adopted in December 2013.\par
\emph{References:} Independent Commission on Banking, \emph{Final Report}, Sept 2011; Dodd--Frank Act, sec.~619.
\end{history}

% ------------------------------------------------------------------
\section{The Eurodollar market in crisis}\label{sec:ed-crisis}

The chart shown plots the Eurodollar rate (LIBOR) minus OIS, the \emph{overnight index swap} rate, ``a sort of term Fed funds'', at one and three months, through the
crisis \ts{8}{2}{0:28}\ts{8}{2}{0:48}\ts{8}{2}{1:08}. Normally the spread is a few basis points, and one expects it: banks can borrow and lend in either market, and any gap is an
arbitrage \ts{8}{2}{1:50}\ts{8}{2}{2:11}. In August 2007 it jumped from almost zero to almost 100 basis points: a measure of the disorganisation of the world money market
\ts{8}{2}{2:31}\ts{8}{2}{2:51}. Mehrling was teaching the course that fall, and told the class something very bad was happening: banks without access to the Fed funds market
needed to roll their one- or three-month funding, and banks with access would not lend to them; the rate rose to offer a huge incentive to borrow Fed funds and lend at LIBOR
\ts{8}{2}{3:12}\ts{8}{2}{3:32}\ts{8}{2}{3:53}. After Bear Stearns the Fed seemed to stabilise the spread, high but steady; it lost control at Lehman, regained it, and the crisis
subsided \ts{8}{2}{3:53}\ts{8}{2}{4:13}.

\begin{definition}[New concepts: basis point, overnight index swap]
A \emph{basis point} is one hundredth of a percentage point (2\% $+$ 1 bp $=$ 2.01\%); in money markets, with sums of millions, even the fourth decimal matters
\ts{8}{2}{1:08}\ts{8}{2}{1:29}. An \emph{overnight index swap} (OIS) exchanges a fixed rate for the average of the overnight rate (in dollars, Fed funds effective) over a
term; its fixed rate is the market's term version of Fed funds. LIBOR $-$ OIS measures the premium for term unsecured interbank funding.
\end{definition}

\begin{nfigure}
\begin{tikzpicture}[font=\scriptsize,x=0.9cm,y=0.9cm]
  \draw[-{Stealth[length=3pt]}] (0,0) -- (9,0) node[right]{time};
  \draw[-{Stealth[length=3pt]}] (0,0) -- (0,4) node[above]{LIBOR $-$ OIS (bp)};
  \foreach \y/\l in {1/90,2/180,3/270} {\draw (-0.08,\y) -- (0.08,\y); \node[left] at (0,\y) {\l};}
  \draw[thick,thmcol] plot[smooth] coordinates {(0,0.05) (1.8,0.06) (2,0.9) (2.4,0.6) (2.8,1.0) (3.2,0.7) (3.6,0.85) (4.2,0.8) (4.8,0.85) (5.2,3.6) (5.6,2.4) (6.2,1.2) (7,0.4) (8.5,0.12)};
  \draw[dashed] (2,0) -- (2,3.8) node[above]{Aug 2007};
  \draw[dashed] (3.6,0) -- (3.6,3.8) node[above]{Bear Stearns};
  \draw[dashed] (5.2,0) -- (5.2,3.8) node[above]{Lehman};
\end{tikzpicture}
\caption{Schematic of the LIBOR--OIS spread through the crisis \ts{8}{2}{0:28}: near zero before August 2007, around 100 basis points after it, steadied after Bear
Stearns, a peak at Lehman, then back to normal. (Shape illustrative; the three-month spread peaked at about 365 bp in October 2008.)}\label{fig:ed-libor}
\end{nfigure}

% ------------------------------------------------------------------
\section{What are Eurodollars?}\label{sec:ed-what}

Stigum's chapter is a recent addition and ``still doesn't go far enough'' \ts{8}{3}{0:07}. Her example (p.~212): Exxon moves money from its account at Chase New
York to Citibank London, a branch outside the country \ts{8}{3}{0:28}\ts{8}{3}{0:49}\ts{8}{3}{1:09}. Behind the scenes Chase transfers reserves at the New York Fed to
Citi New York, which acts as correspondent of Citi London \ts{8}{3}{1:55}\ts{8}{3}{2:15}\ts{8}{3}{2:37}:
\begin{taccts}
\tacct[2.2cm]{Exxon}{$-$Deposit at Chase\\$+$Deposit at Citi London}{}\quad
\tacct[2.2cm]{Chase NY}{$-$Reserves}{$-$Deposit of Exxon}

\medskip
\tacct[2.2cm]{Citi NY}{$+$Reserves}{$+$Deposit of Citi London}\quad
\tacct[2.2cm]{Citi London}{$+$Deposit at Citi NY}{$+$Deposit of Exxon}
\end{taccts}
Stigum stresses that ``no money ever leaves the United States'': reserves at the New York Fed only change owner, from Chase to Citi \ts{8}{3}{3:12}\ts{8}{3}{3:32}. True,
but a \emph{deposit} leaves the country, and that is money; and the deposit of Citi London at Citi New York is, for Citi London, its \emph{reserves}: it cannot hold an
account at the Fed, but it can hold one at a bank that does \ts{8}{3}{3:53}\ts{8}{3}{4:14}\ts{8}{3}{4:34}.

\paragraph{A clearer example.} A branch makes it look like internal accounting of one bank. Better: Exxon moves its money to Crédit Lyonnais in Paris, which holds its
dollar reserves as a deposit at Citi New York, a different bank: a correspondent relationship, and Citi New York certainly considers its reserves to back \emph{its own}
deposits \ts{8}{3}{4:55}\ts{8}{3}{5:16}\ts{8}{3}{5:36}\ts{8}{3}{5:57}. Crédit Lyonnais can now go into the dollar business: make dollar loans by creating dollar deposits, and
pay withdrawals with its deposit in New York, even though it is in France \ts{8}{3}{6:18}\ts{8}{3}{6:38}\ts{8}{3}{6:59}.

\begin{definition}[New concepts: Eurodollar]\label{def:ed}
\emph{Eurodollars} are dollar-denominated deposits (and loans) at banks outside the United States, including foreign branches of US banks. Their reserves are deposits at
banks in the United States, not at the Fed. (``Euro'' because the market started in Europe; the name is used for offshore dollars anywhere, and has nothing to do with the
euro currency.)
\end{definition}\begin{history}[Origins of the Eurodollar market]
The market grew in London in the late 1950s. Commonly cited causes: dollar balances held in Europe by Soviet-bloc banks wary of keeping them in the US (the Paris bank BCEN,
whose telex address was ``Eurobank'', is often credited with the name); the 1957 sterling crisis, when the UK restricted sterling credit to non-residents and London banks
turned to dollars; US Regulation Q ceilings on deposit rates; and later US capital controls (the Interest Equalization Tax of 1963, the voluntary credit restraints of 1965).\par
\emph{References:} C.~R.~Schenk, ``The origins of the Eurodollar market in London: 1955--1963'', \emph{Explorations in Economic History} 35 (1998).
\end{history}

\section{Why is there a Eurodollar market?}\label{sec:ed-why}
There is demand: the dollar is the world reserve currency, and traders in Japan and Argentina may invoice and pay in dollars without any office in the US
\ts{8}{4}{0:07}\ts{8}{4}{0:28}. They could do it through accounts in New York, but the Fed prefers that this happens offshore: it wants to focus on domestic credit
\ts{8}{4}{0:49}. And a payment system is a credit system: deficit agents must be able to borrow from surplus agents, so people want not only dollar deposits but dollar loans
\ts{8}{4}{1:09}\ts{8}{4}{1:30}. A bank like Crédit Lyonnais wants a matched book in dollars, as many dollar assets as dollar liabilities, so that it has no currency exposure in
its euro accounts \ts{8}{4}{1:50}.

\paragraph{History and size.} The market arose because of controls on international capital flows after the Second World War: large corporations did their dollar business
with each other in London, evading US controls, and the US let them, even encouraged them; the controls broke down, then Bretton Woods, and the market kept thriving, growing
an order of magnitude larger than the Fed funds market \ts{8}{4}{2:13}\ts{8}{4}{2:33}\ts{8}{4}{2:54}.

\paragraph{Tail and dog.} Before the crisis it was unclear which was the market for dollars, Fed funds or Eurodollars: they traded at the same rate. People at the New York Fed,
and Mehrling too, worried that with so much volume offshore, and big banks able to borrow in either market, the Fed could lose control of the rate
\ts{8}{4}{3:14}\ts{8}{4}{3:55}\ts{8}{4}{4:15}. The crisis settled it: the Eurodollar market froze, the Fed funds market did not, because the Fed supported it. Fed funds are the
``real'' dollars, Eurodollars a credit extension of them, below Fed funds in the hierarchy \ts{8}{4}{3:34}\ts{8}{4}{4:36}.

% ------------------------------------------------------------------
\section{The Eurodollar market as global funding market}\label{sec:ed-funding}

The Eurodollar market is not only a payment system: it is \emph{the} world funding market, and the world funding market is a dollar market \ts{8}{5}{0:08}\ts{8}{5}{0:29}.
Ultimate lenders and borrowers anywhere want to lend and borrow in dollars, for periods of time, and the Eurodollar system intermediates the flow: an overnight deposit (a money
balance, perhaps petrodollars) funds a six-month term loan, perhaps to a Korean manufacturer \ts{8}{5}{0:50}\ts{8}{5}{1:10}\ts{8}{5}{1:30}\ts{8}{5}{1:54}.
\begin{taccts}
\tacct[2.4cm]{Eurodollar system}{Term loan (6 months)}{Overnight deposit}
\end{taccts}
This is the repo dealer's balance sheet of lecture~7, overnight funding against term lending: here a bank acts as a dealer, quoting a deposit rate and a loan rate, a bid--ask
spread. ``Banks are dealers in money'' \ts{8}{5}{2:14}\ts{8}{5}{2:34}\ts{8}{5}{2:54}. Unlike repo, Eurodollars are \emph{unsecured}: the security is the whole balance sheet of the
bank, its capital, and its backstop, the national government; unlike Fed funds, a brokered market, Eurodollar funding shows up on the balance sheets of the Eurobanks
\ts{8}{5}{3:15}\ts{8}{5}{3:36}\ts{8}{5}{3:57}. A small country borrowing large amounts is most likely funded, ultimately, in the world dollar market; borrowers and depositors outside
the US face exchange risk, a topic for later \ts{8}{5}{4:17}\ts{8}{5}{4:38}\ts{8}{5}{4:58}.

\paragraph{Outside US regulation.} By the home-country rule a bank is supervised where it resides: a French bank dealing in dollars is supervised by the Banque de France and to
some extent the ECB, not by the Fed. It can expand its balance sheet in dollars and create what is essentially money, not counted in the US money supply, since it is neither
the liability of an onshore bank nor the asset of a US resident. Yet for Exxon a deposit in Europe is ``more or less the same thing'' as one in New York: they trade at par
\ts{8}{5}{5:19}\ts{8}{5}{5:39}\ts{8}{5}{5:59}\ts{8}{5}{6:20}. The crisis chart shows the stress on that par: to keep it, Eurodollars had to pay a full percentage point more, to keep
people from ``breaking the buck'' \ts{8}{5}{6:20}\ts{8}{5}{6:41}.

\paragraph{The interbank market.} Individual banks are not matched. Say Crédit Lyonnais has many customers who want dollar deposits, more than it can lend: excess dollar
deposits; Citi London can lend any amount but has a thin European deposit base: excess dollar loans. They do business with each other through an interbank loan
\ts{8}{5}{7:01}\ts{8}{5}{7:22}\ts{8}{5}{8:08}\ts{8}{5}{8:30}\ts{8}{5}{9:13}:
\begin{taccts}
\tacct[2.4cm]{Crédit Lyonnais}{Interbank loan to Citi London}{Dollar deposits}\qquad
\tacct[2.4cm]{Citi London}{Dollar loans}{Interbank deposit from Crédit Lyonnais}
\end{taccts}
LIBOR, the London interbank offered rate, is the rate on this interbank lending \ts{8}{5}{9:47}\ts{8}{5}{10:08}. People ask why banks do all this lending to each other instead of
lending to ``real productive activity'': this is why, to move money from where it is to where it is needed \ts{8}{5}{10:28}\ts{8}{5}{10:50}. It is this market that is breaking down
in Europe now: the loans are unsecured, and a bank that is not sure it can lend on may not want to take the deposit in the first place \ts{8}{5}{11:11}\ts{8}{5}{11:31}.

% ------------------------------------------------------------------
\section{The liquidity challenge of Eurodollar banks}\label{sec:ed-liquidity}

Eurobanks hold their dollar reserves at a New York bank, not at the Fed: they cannot borrow Fed funds, nor borrow from the Fed; their backstop is their own central bank, and the
ECB ``can't print dollars''. So they face more discipline than onshore banks \ts{8}{6}{0:07}\ts{8}{6}{0:27}\ts{8}{6}{0:48}. They respond by lining up their cash flows in time much
more carefully: deposits are \emph{term} deposits to a specific date (one month, three months), loans run to specific dates (six months), and the banks match not only the
amounts but the \emph{dates}, so that when they must pay, someone is paying them \ts{8}{6}{1:09}\ts{8}{6}{1:29}\ts{8}{6}{1:50}\ts{8}{6}{2:10}. Payments do not always come as promised,
so reserves are still needed, but one can try; and this is the motivation for forward rate agreements and foreign exchange agreements: instruments to line up cash inflows and
outflows in time \ts{8}{6}{2:30}\ts{8}{6}{2:50}\ts{8}{6}{3:10}.

\paragraph{A forward deposit as a swap of IOUs.} Bank X expects to make a three-month loan two months from now (a corporate customer has announced its needs) and wants to lock in
the funding now; bank Y expects a deposit inflow in two months and wants to lock in a use for it \ts{8}{6}{4:55}\ts{8}{6}{5:15}\ts{8}{6}{5:36}\ts{8}{6}{5:56}\ts{8}{6}{6:17}. ``All banking is a swap
of IOUs'': X makes a two-month deposit at Y of \$1 million, Y makes a five-month deposit at X of \$1 million \ts{8}{6}{6:37}\ts{8}{6}{7:02}:
\begin{taccts}
\textit{$t=0$:}\quad\tacct[2.4cm]{Bank X}{Deposit at Y (2 months)}{Deposit of Y (5 months)}\quad
\tacct[2.4cm]{Bank Y}{Deposit at X (5 months)}{Deposit of X (2 months)}

\medskip
\textit{$t=2$:}\quad\tacct[2.4cm]{Bank X}{$-$Deposit at Y\\$+$3-month loan}{Deposit of Y (3 months left)}\quad
\tacct[2.4cm]{Bank Y}{Deposit at X (3 months left)}{$-$Deposit of X\\$+$Customer deposit}
\end{taccts}
No cash moves at time 0 \ts{8}{6}{7:23}\ts{8}{6}{7:44}. At two months X's deposit at Y matures: Y pays X, which uses the money to make its loan; Y pays with the deposit inflow it
expected \ts{8}{6}{7:44}\ts{8}{6}{8:04}\ts{8}{6}{8:26}\ts{8}{6}{8:47}. The swap has locked in a flow of funds from Y to X in two months, for three months; Eurodollar deposits exist to
any date, so the timing can be fine-tuned with balance sheet entries alone \ts{8}{6}{9:08}\ts{8}{6}{9:28}. If the two deposits pay the same rate over the last three months, call it
$f$, X has locked in a funding rate and Y an investment rate, even if three-month LIBOR is much lower by then \ts{8}{6}{9:49}\ts{8}{6}{10:12}\ts{8}{6}{10:32}\ts{8}{6}{10:53}.

\section{FRAs as implicit swaps of IOUs}\label{sec:ed-fra}

The literal swap takes up balance sheet space, against which regulators require capital and reserves; so it is done off balance sheet \ts{8}{7}{0:08}. Stigum describes the
evolution \ts{8}{7}{0:28}\ts{8}{7}{0:50}:
\begin{enumerate}
  \item the \emph{forward forward}: Y promises to make a deposit at X at rate $f$ two months from now;
  \item the \emph{forward rate agreement}: Y promises nothing of the sort; X funds its loan at three-month LIBOR when the time comes, and the two banks act as if there were a
        three-month deposit paying $f-\text{LIBOR}$ between them \ts{8}{7}{1:10}\ts{8}{7}{1:31}.
\end{enumerate}

\begin{definition}[New concepts: forward rate agreement]\label{def:ed-fra}
A \emph{forward rate agreement} (FRA) between two banks fixes today a rate $f$ for a future period (here from month 2 to month 5); at the start of the period they settle the
difference between the reference rate then observed (LIBOR) and $f$, on a notional amount: if LIBOR $>f$ the forward lender pays the forward borrower, if LIBOR $<f$ the
forward borrower pays \ts{8}{6}{3:31}\ts{8}{6}{3:52}\ts{8}{7}{1:55}\ts{8}{7}{2:15}.
\end{definition}

Market people describe it as a bet on LIBOR, and it can be used to speculate; but that is not why it exists \ts{8}{6}{4:13}\ts{8}{6}{4:33}. X funds at LIBOR plus or minus the FRA payment,
i.e.\ at $f$ whatever happens: the FRA implements exactly what the swap of IOUs did, locking in funding, without balance sheet space. ``It looks like a side bet on LIBOR, but it's used to
line up payments in time'' \ts{8}{7}{2:37}\ts{8}{7}{2:58}\ts{8}{7}{3:18}.

\begin{proposition}[Derivatives as implicit swaps of IOUs]
All the derivatives of the course, interest rate swaps, foreign exchange swaps, credit default swaps, can be viewed as implicit, behind-the-scenes swaps of IOUs: they are integrated in
the balance sheet structure, and are ``not really off balance sheet''. The method: ask how one would build the same payoff with a swap of IOUs on bank balance sheets; that tells what
the instrument is for \ts{8}{7}{3:18}\ts{8}{7}{3:39}\ts{8}{7}{3:59}\ts{8}{7}{4:20}.
\end{proposition}

% ------------------------------------------------------------------
\section{Forward parity}\label{sec:ed-parity}

\paragraph{Forward interest parity.} Y would like $f$ high, X low; but they do not haggle: $f$ is fixed by an arbitrage condition \ts{8}{8}{0:07}\ts{8}{8}{0:28}. There are market rates
for two-month and five-month Eurodollar deposits, $R_{0,2}$ and $R_{0,5}$. If the swapped deposits were real deposits at market rates, X would have manufactured a three-month funding
rate and Y an investment rate; the only rate both accept is the one implied by the market rates \ts{8}{8}{0:48}\ts{8}{8}{1:12}\ts{8}{8}{1:37}\ts{8}{8}{2:17}\ts{8}{8}{2:37}:
\begin{keyformulas}
\textbf{Forward interest parity} \ts{8}{8}{1:37}\ts{8}{8}{2:58}:
\[
  (1+R_{0,2})(1+f_{2,5}) = 1+R_{0,5}
\]
(with each rate expressed per its own period; in money-market convention $(1+R_{0,2}\tfrac{60}{360})(1+f_{2,5}\tfrac{90}{360}) = 1+R_{0,5}\tfrac{150}{360}$).

\medskip
\textbf{Covered interest parity} \ts{8}{8}{7:52}\ts{8}{8}{8:12}: with $S$ the spot and $F_6$ the six-month forward exchange rate, both in yen per dollar, and $R$, $R^\ast$ the six-month
dollar and yen deposit rates,
\[
  F_6 = S\,\frac{1+R^\ast}{1+R}.
\]
\end{keyformulas}
FRAs use this $f$ \ts{8}{8}{2:58}.

\paragraph{Covered interest parity.} A bank is asked for a yen loan but has only dollar funding, six-month Eurodollars: matched in time, not in currency \ts{8}{8}{2:58}\ts{8}{8}{3:21}\ts{8}{8}{3:42}.
Forward exchange contracts exist for this; but with a swap of IOUs, it swaps a six-month dollar deposit for a six-month yen deposit with another bank, at today's spot rate
\ts{8}{8}{4:02}\ts{8}{8}{4:22}\ts{8}{8}{4:43}\ts{8}{8}{5:05}. For one dollar: it receives $S$ yen (say 100) to lend; after six months it owes the counterparty $S(1+R^\ast)$ yen and is owed
$1+R$ dollars \ts{8}{8}{5:25}\ts{8}{8}{6:28}\ts{8}{8}{6:48}\ts{8}{8}{7:08}\ts{8}{8}{7:31}. Paying yen and receiving dollars at a ratio fixed today is an implicit forward exchange rate, which must equal the
market forward $F_6$: covered interest parity \ts{8}{8}{7:52}\ts{8}{8}{8:12}. Both parities are almost definitions: since anyone can ``roll their own'' forward from existing instruments, the outright
forward must have the same price, or there is an arbitrage \ts{8}{8}{8:33}\ts{8}{8}{8:53}. (In the yen--dollar example the dollar has no star, ``because the dollar is special'', the world funding
currency \ts{8}{8}{5:46}.)

\begin{taccts}
\tacct[2.4cm]{Bank (yen lender)}{Yen loan $S$\\Dollar deposit at B 1}{Dollar Eurodollar funding 1\\Yen deposit of B $S$}\qquad
\tacct[2.4cm]{Bank B}{Yen deposit at bank $S$}{Dollar deposit of bank 1}
\end{taccts}

\section{Forward rates are not expected spot rates}\label{sec:ed-eh}

Two natural guesses \ts{8}{9}{0:08}\ts{8}{9}{0:29}\ts{8}{9}{0:49}\ts{8}{9}{1:31}:
\begin{itemize}
  \item the forward interest rate $f_{2,5}$ is an unbiased forecast of three-month LIBOR two months from now: the \emph{expectations hypothesis} of the term structure
        (\cref{def:nh-eh}), $f = \E[\text{future spot rate}]$;
  \item the forward exchange rate $F_6$ is an unbiased forecast of the spot rate six months from now: \emph{uncovered interest parity} (UIP), $F_6 = \E[S_6]$, equivalently
        $\E[S_6]/S = (1+R^\ast)/(1+R)$.
\end{itemize}
Both fail empirically, all the time: a puzzle for economics \ts{8}{9}{1:10}\ts{8}{9}{1:52}. Mehrling's answer, to be developed, is that it has to do with the hierarchy of money and credit, with
liquidity: payments today are worth more than promised payments in the future, and not only because of time preference, in foreign exchange as in the term structure \ts{8}{9}{2:12}\ts{8}{9}{2:32}.

\begin{coursenote}[Segment 8.10]
The video ``Forward Rates are NOT Expected Spot Rates'' (\seg{8}{10}{L8.10}), filed in the playlist under lecture~21, repeats almost word for word the second half of L8.9.
\end{coursenote}

% ------------------------------------------------------------------
\flushside
\section*{Key formulas and ideas}
\begin{keyformulas}
\begin{tabularx}{\linewidth}{@{}l L@{}}
Eurodollars & dollar deposits outside the US; reserves $=$ deposits at a New York correspondent; below Fed funds in the hierarchy\\
Global funding & Eurobanks as money dealers: overnight/term deposits fund term loans; LIBOR $=$ interbank rate\\
Volcker rule & matched book (market making) allowed, proprietary positions not; basis risk blurs the line\\
Liquidity & no Fed access $\Rightarrow$ match amounts \emph{and} dates of cash flows\\
FRA & implicit swap of term deposits; locks in a future funding/investment rate $f$ off balance sheet\\
Forward interest parity & $(1+R_{0,2})(1+f_{2,5}) = 1+R_{0,5}$\\
Covered interest parity & $F = S(1+R^\ast)/(1+R)$\\
EH, UIP & $f=\E[\text{future rate}]$, $F=\E[S_T]$: both fail empirically (liquidity)\\
\end{tabularx}
\end{keyformulas}

\section*{Exercises}

\begin{exercise}[Forward interest parity]
Two-month Eurodollar deposits pay 0.30\% and five-month deposits 0.45\% (annualised, actual/360, 60 and 150 days). Compute the implied three-month forward rate two months ahead. If an
FRA is quoted at 0.70\%, describe the arbitrage.
\end{exercise}

\begin{exercise}[FRA settlement]
Bank X buys an FRA at $f=0.60\%$ on \$10 million for the period from month 2 to month 5 (90 days). At month 2 three-month LIBOR is 0.90\%.
\begin{enumerate}[(a)]
  \item Who pays whom, and how much (paid at the end of the period)?
  \item Show that X's all-in funding cost for its loan is 0.60\%.
\end{enumerate}
\end{exercise}

\begin{exercise}[Covered interest parity]
The spot rate is 80 yen per dollar, six-month dollar deposits pay 0.5\% and yen deposits 0.1\% (for six months).
\begin{enumerate}[(a)]
  \item Compute the six-month forward rate.
  \item Under UIP, what is the expected spot rate in six months? What does the failure of UIP mean for a bank that funds in yen and lends in dollars?
\end{enumerate}
\end{exercise}

\begin{exercise}[Exxon's transfer]
Redo Stigum's example with Crédit Lyonnais in place of Citi London, and then let Crédit Lyonnais lend \$5 million to a Korean firm that pays a US supplier banking at Chase. Write all the
T-accounts, including the Fed.
\end{exercise}
